\section{流形是局部像欧氏空间的集合}
\label{sec:12-Real-Analysis-Support/08-Manifold-Basics/01}

你用过地图。打开手机地图，屏幕上的平面图像如实反映了你周围几百米的街道布局——转角的角度、道路的相对位置，都和真实世界一致。但你知道没有一张平面地图能同时画下整个地球而不失真。墨卡托投影把格陵兰画得和非洲一样大；任何平面世界地图都会在某些地方扭曲形状、面积或角度。

这不是画图技术的问题。球面和平面在拓扑上根本不同胚——不存在一个连续且双向连续的一一映射，能把整个球面平铺到平面上，同时保持所有几何关系。数学不做违反拓扑的事。

但一本地图册可以解决问题。每页负责一块区域，翻页处重叠的部分彼此对得上。流形定义的正是这样一个对象：整体可以弯曲甚至复杂，局部却像普通的欧氏空间。

\subsection{坐标卡：把一小块拉平}

设 \(M\) 是一个拓扑空间——先不要求它是任何欧氏空间的子集，只需它有开集、连续这些基本概念。一个 \(d\) 维坐标卡（chart）是一对 \((U,\varphi)\)，其中 \(U\subseteq M\) 是开集，

\[
\varphi:U\to\varphi(U)\subseteq\R^d
\]

是把 \(U\) 同胚地映到 \(\R^d\) 中开集的映射。

"同胚"意味着双向连续的一一对应。翻译成操作语言：坐标的读数和点本身可以互相还原——给定点，你能算出坐标；给定坐标，你能找回那个点。并且接近的点对应接近的坐标，不会出现两个靠得很近的点在地图坐标上突然跳开。反过来，地图上相邻的两个坐标值，对应原始空间中相邻的两个点。

一张卡相当于一页地图：给你一块区域和一套坐标系。但通常一张卡不够——你需要一整本图册。

\subsection{图册与转移映射：翻页处的对照表}

一组覆盖整个 \(M\) 的坐标卡称为图册（atlas）。这里的关键词是"覆盖"：\(M\) 的每个点至少要落在一张卡的范围内。如果有两张卡 \((U,\varphi)\) 和 \((V,\psi)\) 覆盖了同一块区域，重叠部分 \(U\cap V\) 中的每个点在两卡中各有坐标。转移映射

\[
\psi\circ\varphi^{-1}:
\varphi(U\cap V)\to\psi(U\cap V)
\]

告诉你同一组点怎样从一页地图的坐标换算到另一页地图的坐标。画成图就是：点在 \(U\) 上的坐标 \(\varphi(p)=u\)，换算到 \(V\) 上的坐标 \(\psi(p)=\psi(\varphi^{-1}(u))\)。

如果图册中任意两张重叠的卡的转移映射都光滑（无限次可微），这个图册就把 \(M\) 变成了一个 \(d\) 维光滑流形。

把光滑性要求放在转移映射上，而不是直接放在 \(\varphi\) 上，这个选择有深层原因。\(M\) 上根本没有全局坐标，"光滑"不能直接定义——它需要一个 \(\R^d\) 内部的参照系。转移映射在 \(\R^d\) 的两个开集之间来回，是标准的多元微积分对象，光滑性有明确定义。因此流形的微分结构就借由这些"翻页操作"来定义。

维数 \(d\) 是良定义的：同一点附近任意两张卡的转移映射都定义在 \(\R^d\) 的开子集之间，维数必须一致，否则不存在光滑可逆的映射。你不能有一张卡把一块区域映成二维的，另一张卡把同一块区域映成三维的——这两张卡的转移映射在 2 维和 3 维之间无法光滑可逆。

\subsection{二维球面：最小非平凡例子}

二维球面

\[
S^2=\{x\in\R^3:\norm{x}=1\}
\]

是最简单的非平凡光滑流形。说它非平凡，是因为它不能用一张卡拉平——\(S^2\) 是紧的（有界闭），而任何单张卡的像都是 \(\R^2\) 的开集，紧集不可能同胚于开集。拓扑本身封死了单卡的可能。

但两张卡就够。从北极 \((0,0,1)\) 做球极投影：对北极以外的任意一点 \(p\in S^2\)，连接北极和 \(p\) 的直线与赤道平面 \(z=0\) 交于一点，把这个交点的二维坐标赋给 \(p\)。这张卡覆盖了"球面去掉北极"。同样从南极再做一张卡，覆盖"球面去掉南极"。两卡在球面去掉两极处重叠——赤道附近的大部分区域同时落在两卡之中。

重叠处的转移映射是什么？球极投影的公式是 \(y = (x_1,x_2)/(1-x_3)\)，逆映射能写出来。两张卡之间的转移映射最后化简为 \(\R^2\setminus\{0\}\) 上的反演 \(y\mapsto y/\norm y^2\)，它是光滑的。两张卡合起来，\(S^2\) 成为一个二维光滑流形。

"地球仪与地图册"的类比在这里完全精确。每张地图是坐标卡，地图册是图册，翻页处的比例换算就是转移映射。失真不是因为地图画得差，而是因为 \(S^2\) 的整体弯曲无法被单张卡消除；流形定义接受这个现实，只要求逐页可用。

\subsection{三个反面教材：什么不是流形}

流形的定义是一组精确的条件。看看什么东西不满足这些条件，边界就更清楚了。

\begin{itemize}
\tightlist
\item
  \textbf{圆锥顶点。} 圆锥面在顶点附近的任何邻域都无法同胚于平面开集——顶点是一个"尖点"，无论你怎么拉平，它都会在映射中暴露出非光滑的褶皱。围绕顶点的环路在锥面上可以缩成一段短弧，但在平面上去掉一点后的小环路不能连续缩成一点。圆锥面去掉顶点是流形，带上顶点不是。
\item
  \textbf{两条相交直线。} 在 \(\R^2\) 中取两条过原点的直线。交点附近的任何小邻域都像"十"字，而一维流形的局部模型是一段线段。"十"字不能同胚于一段线段——除去交点后，邻域分裂成四个连通分支，而线段除去一点后只分裂成两个。这不是一维流形。
\item
  \textbf{带边区间 \([0,1)\)。} 端点 \(0\) 的任何小邻域在区间内只与半开区间 \([0,\varepsilon)\) 同胚，不与 \(\R\) 的开集同胚（开集的同胚像仍是开的）。它是带边流形，属于稍宽的定义，但本书默认的"流形"不含边。
\end{itemize}

判断一个集合是不是流形，要问的是：每一点附近能否找到一张合法的坐标卡。整体形状——是否闭合、是否打结、有几个洞——不影响这个逐点的局部性质，那些是整体拓扑另外关心的事。

\subsection{子流形：生活在更大的空间里}

本书遇到的流形绝大多数不是抽象的拓扑空间，而是老老实实地嵌在某个 \(\R^n\) 里面。这类流形叫子流形，有更直接的判断方法。

若 \(M\subseteq\R^n\)，且对 \(M\) 上的每一点，都存在一个邻域和 \(\R^n\) 上的局部坐标变换，使 \(M\) 在这个新坐标系下恰好对应其中某 \(d\) 个坐标自由变化、其余 \(n-d\) 个坐标固定为零，则 \(M\) 是 \(\R^n\) 的 \(d\) 维子流形。球面、环面、以及 Part 13 中数据集中分布在上的低维曲面，都是这个意义上的对象。

最实用的判据是正则值定理：若 \(F:\R^n\to\R^k\) 光滑，且对某个 \(c\in\R^k\)，在 \(F^{-1}(c)\) 上的每一点，\(F\) 的 Jacobi 矩阵都满秩（秩为 \(k\)），则 \(F^{-1}(c)\) 是 \(\R^n\) 的 \(n-k\) 维子流形。球面恰好是这个形式的特例：取 \(F(x)=\norm x^2\) 和 \(c=1\)，则 \(\nabla F = 2x \neq 0\) 在球面上处处成立（球面不包含原点），Jacobi 满秩，因此球面是 \(n-1\) 维子流形。

这个定理不需要积分、不需要概率，只需要微积分中 Jacobi 满秩的条件，就能从方程约束反推几何结构。

\subsection{为什么要跑到 Part 13 才用这个定义}

在机器学习和深度学习中，"流形假设"经常被挂在嘴边：高维观测集中在低维流形附近。有了流形的正式定义，这句话的三层含义可以清楚地分开，不再混在一起当口号用。

第一层：\textbf{数据有结构。} 观测并不充满整个环境空间——比如 784 维空间中手写数字的分布只占据极低维的区域。这是经验事实，检查数据就能验证。

第二层：\textbf{结构是流形。} "局部可拉平、有确定统一的维数、没有尖点或自交"——这是建模假设，不是数据天然具备的性质。数据完全可能躺在圆锥顶点上，或分布在两条相交的低维平面上。那些情形下"维数"和"近邻"的概念都需要重新解释，标准的流形方法不能直接套用。

第三层：\textbf{局部是平的。} 每个点附近存在近似的切空间——这是光滑流形的定理，不是额外假设。只要第二层成立，这一层自动成立。下一节\hyperref[sec:12-Real-Analysis-Support/08-Manifold-Basics/02]{``切空间是流形的一阶线性近似''}会把这个定理变成切空间的正式定义。

三层分清之后，你才能判断：说某个数据集"满足流形假设"时，你验证了什么（数据不充满全空间），假设了什么（局部没有奇点），以及可以推导出什么（存在切空间结构）。
